% Opening of Chapter 1 (Introduction)

The Yersinia pestis bacterium is the causative agent of bubonic plague. Perhaps more than any other human pathogen, it has struck fear into the hearts and minds of peoples all over the world, and consciously or not lingers as a shadow in the western psyche. 

The present thesis will explore the mechanisms this bacteria employs in a human host, and the strategies it has evolved to overcome our innate immune system which possibly facilitated its historical success. 

We will paint a mathematical picture of an early stage plague infection with the aim of shedding light on the challenges faced by Yersinia pestis, and potentially the strategic advantege of its niche method.
